\documentclass[a4paper,10pt]{article}

\usepackage[utf8]{inputenc}
\usepackage[T1]{fontenc}
\usepackage{geometry}
\usepackage{booktabs}
\usepackage{array}
\usepackage{tabularx}
\usepackage{longtable}
\usepackage{listings}
\usepackage{xcolor}
\usepackage[hidelinks]{hyperref}
\usepackage{amsmath}

\geometry{
    a4paper,
    top=1.8cm,
    bottom=1.8cm,
    left=1.8cm,
    right=1.8cm
}

\setlength{\parindent}{0pt}
\setlength{\parskip}{5pt}
\renewcommand{\arraystretch}{1.18}

\lstdefinestyle{cstyle}{
    language={[ANSI]C},
    basicstyle=\ttfamily\small,
    keywordstyle=\bfseries,
    commentstyle=\itshape,
    numbers=left,
    numberstyle=\tiny,
    numbersep=7pt,
    frame=single,
    breaklines=true,
    showstringspaces=false,
    tabsize=4
}

\newcommand{\hex}[1]{\texttt{0x#1}}
\newcommand{\code}[1]{\texttt{#1}}

\title{\textbf{DEOS Lighting Service ICD}\\
\large Sürüm 1.2 -- Aydınlatma Servisi (Hedef: Main STM32)}
\author{\textbf{Emir Arkalı} (\href{mailto:emirarkali07@gmail.com}{emirarkali07@gmail.com})\\DEOS -- Dokuz Eylül Otonom Sistemler}
\date{İlk Yayın: 28 Eylül 2026 \quad | \quad Son Revizyon: 14 Ekim 2026}

\begin{document}
\maketitle

\begin{center}
\small
\begin{tabularx}{\textwidth}{ll lX}
\toprule
\textbf{Doküman Tipi:} & Interface Control Document (ICD) & \textbf{Doküman Kodu:} & \code{DOC-ICD-SRV-LGHT-012} \\
\textbf{Doküman Sürümü:} & Sürüm 1.2 & \textbf{Protokol Sürümü:} & v1.1 (\hex{11}) -- Medya Bağımsız \\
\textbf{Yazar:} & Emir Arkalı (\texttt{emirarkali07@gmail.com}) & \textbf{Son Değiştiren:} & Emir Arkalı \\
\textbf{İlk Yayın Tarihi:} & 28 Eylül 2026 & \textbf{Son Revizyon Tarihi:} & 14 Ekim 2026 \\
\textbf{Durum:} & Yürürlükte / Onaylandı & \textbf{Taşıma / Kapsam:} & LIGHTING Servisi (\hex{08}) -- Hedef: Main STM32 (\hex{03}) \\
\bottomrule
\end{tabularx}
\end{center}

\subsection*{Doküman Revizyon Geçmişi}
\begin{center}
\small
\begin{tabularx}{\textwidth}{c c l X}
\toprule
\textbf{Revizyon} & \textbf{Tarih} & \textbf{Yazar / Değiştiren} & \textbf{Açıklama} \\
\midrule
v0.3 & 28 Eylül 2026 & Emir Arkalı & İlk taslak ve monolitik ICD v0.3 tanımı. \\
v1.0 & 04 Ekim 2026 & Emir Arkalı & Aydınlatma alt sistem ICD'si oluşturuldu. \\
v1.1 & 13 Ekim 2026 & Emir Arkalı & Sürüm 1.1 standardizasyonu, servisin Main STM32 üzerinde çalıştığı netleştirildi, kurumsal üstveri eklendi. \\
v1.2 & 14 Ekim 2026 & Emir Arkalı & Durum telemetrisi (\code{GET\_STATUS} / \code{STATUS}) ve geri bildirim payload yapısı eklendi. \\
\bottomrule
\end{tabularx}
\end{center}
\vspace{0.8em}

\section{Aydınlatma (Lighting) Servisi Mimarisi}

DEOS dağıtık mimarisinde Aydınlatma (Lighting) bağımsız bir fiziksel mikrodenetleyici veya mantıksal düğüm (Node) değildir; araç ana kontrol beyni olan \code{DEOS\_NODE\_MAIN\_STM32} (\hex{03}) üzerinde doğrudan koşan fonksiyonel bir servistir (\code{DEOS\_SERVICE\_LIGHTING}, \hex{08}).

Dolayısıyla, yer istasyonundan (Ground Control), Textual arayüzden veya otonom planlayıcıdan gönderilecek tüm aydınlatma mesajlarının hedef düğüm adresi (Destination Node) her zaman \textbf{Main STM32} olmalıdır:

\begin{itemize}
    \item \textbf{Hedef Düğüm (Destination Node):} \code{DEOS\_NODE\_MAIN\_STM32} (\hex{03})
    \item \textbf{Fonksiyonel Servis (Service ID):} \code{DEOS\_SERVICE\_LIGHTING} (\hex{08})
    \item \textbf{Mesaj Sınıfı (Message Class):} \code{DEOS\_CLASS\_COMMAND} (\hex{01}), \code{DEOS\_CLASS\_REQUEST} (\hex{04}), \code{DEOS\_CLASS\_STATUS} (\hex{02}), \code{DEOS\_CLASS\_RESPONSE} (\hex{05})
\end{itemize}

\subsection{Far ID Tanımları (Light ID)}

Komut gönderilirken hangi farın hedeflendiği 1 baytlık \code{Light ID} ile belirtilir:

\begin{center}
\begin{tabularx}{\textwidth}{c l l X}
\toprule
\textbf{Far ID} & \textbf{Sembolik İsim} & \textbf{Far / Lamba} & \textbf{Açıklama} \\
\midrule
\hex{01} & \code{DEOS\_LIGHT\_HEADLIGHT} & Ön Farlar & Gece sürüşü ana aydınlatma. \\
\hex{02} & \code{DEOS\_LIGHT\_BRAKE} & Fren Lambası & Arka stop ikaz lambası. \\
\hex{03} & \code{DEOS\_LIGHT\_TURN\_LEFT} & Sol Sinyal & Sol dönüş lambası. \\
\hex{04} & \code{DEOS\_LIGHT\_TURN\_RIGHT} & Sağ Sinyal & Sağ dönüş lambası. \\
\hex{FF} & \code{DEOS\_LIGHT\_ALL} & Tüm Işıklar & Tüm far grubunu kapsayan genel ID. \\
\bottomrule
\end{tabularx}
\end{center}

\subsection{Aydınlatma Komutları ve Telemetri}

Aydınlatma servisi kontrol komutlarını ve durum sorgulama/yayınlama mekanizmasını destekler:

\begin{center}
\begin{tabularx}{\textwidth}{c l c c X}
\toprule
\textbf{Komut ID} & \textbf{Sembolik İsim} & \textbf{Sınıf} & \textbf{Length} & \textbf{Payload İçeriği} \\
\midrule
\hex{01} & \code{DEOS\_CMD\_LIGHTING\_TURN\_ON} & COMMAND & 1 byte & Hedef \code{Light ID} (Açılacak far) \\
\hex{02} & \code{DEOS\_CMD\_LIGHTING\_TURN\_OFF} & COMMAND & 1 byte & Hedef \code{Light ID} (Kapatılacak far) \\
\hex{03} & \code{DEOS\_CMD\_LIGHTING\_GET\_STATUS} & REQUEST & 0--1 byte & Opsiyonel \code{Light ID} (\hex{FF} veya boş: tümü) \\
\hex{04} & \code{DEOS\_CMD\_LIGHTING\_STATUS} & STATUS/RESP & 6 byte & Aydınlatma durum yapısı (\code{deos\_lighting\_status\_payload}) \\
\bottomrule
\end{tabularx}
\end{center}

\subsection{Komut Payload Yerleşimi (TURN\_ON / TURN\_OFF)}

Açma/kapama komutlarında DEOS ortak başlığından sonra gelen payload alanı 1 bayttır:

\begin{center}
\begin{tabular}{c c c l}
\toprule
\textbf{Offset} & \textbf{Boyut} & \textbf{Tip} & \textbf{Alan Tanımı} \\
\midrule
0 & 1 byte & \code{uint8\_t} & \textbf{Light ID}: \hex{01} Ön Far, \hex{02} Fren, \hex{03} Sol Sinyal, \hex{04} Sağ Sinyal, \hex{FF} Tümü \\
\bottomrule
\end{tabular}
\end{center}

\subsection{Aydınlatma Durum Telemetrisi (STATUS / RESPONSE)}

\code{DEOS\_CMD\_LIGHTING\_STATUS} (\hex{04}) paketi, otonomi katmanı veya yer istasyonunun talebi (\code{GET\_STATUS}) üzerine cevap olarak döner; ayrıca lamba durumu değiştiğinde anlık (event-triggered) ve periyodik olarak 1~Hz frekansla CAN Broadcast (\hex{FE}) üzerinden yayınlanır:

\begin{center}
\begin{tabular}{c c c l}
\toprule
\textbf{Offset} & \textbf{Boyut} & \textbf{Tip} & \textbf{Alan Tanımı} \\
\midrule
0 & 1 byte & \code{uint8\_t} & \textbf{headlight\_state}: \hex{00} OFF, \hex{01} ON, \hex{02} BLINKING, \hex{03} FAULT \\
1 & 1 byte & \code{uint8\_t} & \textbf{brake\_light\_state}: \hex{00} OFF, \hex{01} ON, \hex{02} BLINKING, \hex{03} FAULT \\
2 & 1 byte & \code{uint8\_t} & \textbf{turn\_left\_state}: \hex{00} OFF, \hex{01} ON, \hex{02} BLINKING, \hex{03} FAULT \\
3 & 1 byte & \code{uint8\_t} & \textbf{turn\_right\_state}: \hex{00} OFF, \hex{01} ON, \hex{02} BLINKING, \hex{03} FAULT \\
4 & 1 byte & \code{uint8\_t} & \textbf{hazard\_state}: \hex{00} OFF, \hex{01} ON, \hex{02} BLINKING, \hex{03} FAULT \\
5 & 1 byte & \code{uint8\_t} & \textbf{fault\_flags}: Arıza bitmask'i (b0: Ön Far, b1: Fren, b2: Sol, b3: Sağ) \\
\bottomrule
\end{tabular}
\end{center}

\textbf{Toplam Durum Payload Uzunluğu (Length):} 6 byte (\hex{06}).

\section{C/C++ Protokol Tanımları}

\begin{lstlisting}[style=cstyle]
/* Far ID Tanimlari */
typedef enum {
    DEOS_LIGHT_HEADLIGHT  = 0x01, /* On Farlar */
    DEOS_LIGHT_BRAKE      = 0x02, /* Fren Lambasi */
    DEOS_LIGHT_TURN_LEFT  = 0x03, /* Sol Sinyal */
    DEOS_LIGHT_TURN_RIGHT = 0x04, /* Sag Sinyal */
    DEOS_LIGHT_ALL        = 0xFF  /* Tum Isiklar */
} deos_light_id_t;

/* Aydinlatma Komutlari */
typedef enum {
    DEOS_CMD_LIGHTING_TURN_ON    = 0x01, /* Far Ac */
    DEOS_CMD_LIGHTING_TURN_OFF   = 0x02, /* Far Kapat */
    DEOS_CMD_LIGHTING_GET_STATUS = 0x03, /* Durum Sorgula (REQUEST) */
    DEOS_CMD_LIGHTING_STATUS     = 0x04  /* Durum Telemetrisi (STATUS/RESPONSE) */
} deos_lighting_command_t;

/* Aydinlatma Durum Kodlari */
typedef enum {
    DEOS_LIGHT_STATE_OFF      = 0x00, /* Kapali */
    DEOS_LIGHT_STATE_ON       = 0x01, /* Acik (Surekli) */
    DEOS_LIGHT_STATE_BLINKING = 0x02, /* Yanip sonuyor (Flasor) */
    DEOS_LIGHT_STATE_FAULT    = 0x03  /* Ariza (Acik/kisa devre) */
} deos_light_state_t;

/* Aydinlatma Durum Payload Yapisi (6 Byte) */
struct deos_lighting_status_payload {
    uint8_t headlight_state;   /* deos_light_state_t */
    uint8_t brake_light_state; /* deos_light_state_t */
    uint8_t turn_left_state;   /* deos_light_state_t */
    uint8_t turn_right_state;  /* deos_light_state_t */
    uint8_t hazard_state;      /* deos_light_state_t */
    uint8_t fault_flags;       /* Bitmask */
} __attribute__((packed));
\end{lstlisting}

\end{document}
